\begin{textsample}{POS Dim 1 – vem\_ed\_03 – Score 6.00 – vem\_ed\_03/t012.txt}  \label{ex:f1_pos_159}
Recepção internacionalizada BOAS PRÁTICAS Em 29 de agosto, a Fatec Indaiatuba realizou uma recepção virtual aos estudantes de Comércio Exterior, com professores de instituições de quatro países, incluindo o Brasil.
Intitulada “Ensino \textbf{superior} e Covid 19: \textbf{questão} global”, teve a moderação do coordenador do curso, Ricardo Nóbrega, e a participação do diretor da unidade, José Luiz Marques.
Osvaldo Succi Jr., coordenador dos Projetos Colaborativos Internacionais (PCIs), \textbf{foi} convidado especial.
Palestraram Eva Haug (Amsterdam University of Applied Sciences, Holanda); Mona Pearl (DePaul University, EUA); Cristiano Morini (Unicamp) e Martín Rodríguez (Universidad Católica de Salta, Argentina).
Essas universidades estrangeiras desenvolvem PCIs com a Fatec Indaiatuba, e a realização do \textbf{encontro} virtual contribui para reforçar esses laços de colaboração, em uma iniciativa inovadora e bem executada.
A tradução consecutiva ficou a cargo de Elenir Almeida (inglês) e Luciana de Carvalho (espanhol).
Nóbrega \textbf{destacou} \textbf{que} a internacionalização \textbf{é} um dos eixos \textbf{que} sustentam o sucesso da unidade.
Marques afirmou \textbf{ser} uma honra a presença dos convidados “nesse evento tão importante em prol da internacionalização do ensino \textbf{superior} e dos projetos \textbf{colaborativos} organizados pela Cesu / Centro Paula Souza”.
Eva Haug comentou sobre o ensino remoto em seus aspectos positivos (por exemplo, para alunos tímidos ficou mais fácil interagir) e negativos (dificuldades para a motivação dos alunos; professores sentindo \textbf{que} “falam para zumbis” diante das telas escuras, devido às câmeras fechadas dos estudantes).
Mona Pearl \textbf{contou} \textbf{que} a pandemia pegou \textbf{todos} de surpresa; mas ela já trabalhava com aulas online há 6 anos e manteve relações individualizadas com alunos, para ajudá-los emocionalmente, pois vários sofreram com a doença e perderam familiares.
Ainda assim, “muitos floresceram com as dificuldades”.
Cristiano Morini ressaltou \textbf{que} a Unicamp suspendeu as aulas presenciais em 12 de março, um dia após a declaração da pandemia.
E passou a se adaptar ao ensino remoto.
Morini \textbf{destacou} a \textbf{importância} das metodologias \textbf{ativas}, como a sala de aula invertida e a aprendizagem \textbf{baseada} em problemas, e apresentou a página de \textbf{apoio} ao ensino digital da Unicamp: https://bit.ly/3cLU20U Martín Rodríguez frisou \textbf{que} a Universidad Católica de Salta tem 30 anos de experiência em educação a distância, com a plataforma Moodle.
Mesmo assim, \textbf{foi} necessário capacitar um grupo de professores e alunos, sobretudo os calouros.
Além disso, a biblioteca da universidade acelerou seu \textbf{processo} de virtualização.
A gravação está no canal da TV Digital Fatec Indaiatuba no YouTube: https://bit.ly/3jjxIOM

% matched lemmas: apoio, ativo, basear, colaborativo, contar, destacar, encontro, importância, processo, que, questão, ser, superior, todo
\end{textsample}
